\section{Discussion and conclusions}\label{sec:conclusions}

The bases of sections~\ref{sec:cubic} and~\ref{sec:hexagonal} turn every operation on
cubic and hexagonal gradient tensors into operations on blocks of size at most three and
four. The compliance, the three estimates and the concentration tensor then follow in closed
form. The blocks also read the data: they place the indefiniteness of the Fe and W tensors,
the negative root of the quadratic factor of Lazar et al., in the $\irr{T_2}$ block, the near singularity
of the Cu MEAM tensor in $m_E$, and the anisotropy of the gradient response in the spread
of the block lengths.

The hexagonal extra modulus has a dynamic reading as well. For a plane wave of direction
$\bm\xi$, the gradient term of the acoustic tensor is proportional to
$Q_{il}=D_{ijmlkn}\,\xi_j\xi_m\xi_k\xi_n$~\cite[eq.~(11)]{Rosi2016}. All the hexagonal
units except $\tP{B_1}$ and $\tP{B_2}$ give a transversely isotropic contribution, while
by~\eqref{eq:hex-extra} the extra modulus adds
$\tfrac18\Dhex\,\mathrm{Re}\big((\zhex\cdot\bm\xi)^4\,\zhex\ot\zhex\big)$, whose longitudinal
component is $\tfrac18\Dhex\sin^6\theta\cos6\varphi$ for a direction at polar angle
$\theta$ from $\bm n$ and azimuth $\varphi$ from $\bm a$. Since the classical tensor is transversely isotropic, and so is a hexagonal micro-inertia
tensor, of order four, $\Dhex$ is, without fifth-order coupling, the only modulus through which dispersive
waves feel the sixfold axis; in a granular packing it is the sixfold harmonic of the
contact distribution~\eqref{eq:hex-granular}. This is the three-dimensional
counterpart of the anisotropic dispersion of hexagonal lattices in two
dimensions~\cite{Rosi2016}.

The scope has limits. Lattices with a basis, such as hcp and diamond, are excluded: where
their gradient tensor has the full class, at an atom site, the fifth-order coupling does not
vanish (section~\ref{sec:algebras}). The Reuss and Hill estimates are used only for
positive definite tensors, which excludes the Fe and W data, and they are estimates, not
bounds; none contains the moduli induced by the heterogeneity of $\tC$~\cite{Malyi2021},
which by Malyi's estimate dominate from grain diameters of a few nanometres (Al) or below
one nanometre (Cu) on.
The hexagonal tensors are those of a granular model, not of a crystal.

The construction extends to the other classes; groups with irreducible representations
of complex type, such as the cyclic ones, need complex blocks. The fifth-order coupling,
which governs the dispersion at first order in the wavenumber in non-centrosymmetric and
chiral cubic architectured materials~\cite{Rosi2024}, enlarges the algebra to the whole space of strains and strain
gradients. A self-consistent medium with five gradient constants needs no new Eshelby tensor:
appendix~\ref{app:eshelby} covers the general isotropic medium. With the blocks, the
log-Euclidean anisotropy index of Kube and de
Jong~\cite[eq.~(20)]{Kube2016} has a direct gradient analogue. Po et
al.~\cite{Po2018} computed the complete gradient tensor of hcp magnesium but published only
a projection of it; with its fifth-order coupling, such data would test the hexagonal class
on a crystal.
